\documentclass[12pt]{article} 

\usepackage[top=1.5cm, bottom=1.5cm, left=1.5cm, right=1.5cm]{geometry}
\usepackage{color,graphicx,amsmath,mathtools,amssymb,amsthm}
\pagestyle{plain}

\begin{document}  
\pagestyle{empty}
 




%\vskip0.5in


%\noindent {\small This document is property of Northeastern University.  Unauthorized distribution of the content is not~permitted.} \\




\begin{enumerate}



\item  
Let $f(n) = \sum^n_{y{=}1} (n^2 \cdot y^{13})$  \\\\

%Find a simple $g(n)$ such that $f(n) {=} \Theta(g(n))$, by proving that 
%$f(n){=}O(g(n))$, and that $f(n) {=} \Omega(g(n))$.
%
%Solve this using the technique taught in the big-O section of the course. \\
%Don't skip or combine steps. Show every step on a separate line.   
%
% The derivation should be algebraic.  Figures aren't needed. 
%
% Don't expand sums.   Instead use summation notation. 
% 
% 
% If your proof is displayed line by line, do this without redundancy: don't repeat the left side of an inequality or equality, if all you're manipulating is the right side.  (Instead start each new line with the symbol for  inequality or equality)
% 
% Don't use the inequality symbol when you actually have equality.  
% 
% Don't include the problem statement in your submission (except for the first line).  \\\\

{\color{blue}
{\bf Answer:}


}






%%%%%%%%%%%%%%%%%%%%%%%%%%%%%%%%%%%%%%%%%
\newpage

\item
Let $n\geqslant 3$.
Consider an ordered list of numbers, $L = (L_1,L_2,\ldots,L_{n-1},L_n)$

$\Large X_k = \sum\limits_{i=2}^k (L_i - L_{i-1}) $ 

$\Large Y = \sum\limits_{j=3}^n (X_j - X_{j-1}) $ 

%Derive a simple expression for $Y$, in terms of specific values from $L$. 
%The derivation should be stated efficiently and the final answer should be as simple as   possible.  \\



{\color{blue}
{\bf Answer:}


}



\end{enumerate} 




\end{document}






















